\documentclass{article}
\usepackage[T1]{fontenc}
\usepackage{lmodern}
\usepackage{microtype}
\setlength{\emergencystretch}{1em}
\usepackage[polish]{babel}
\usepackage{amssymb}
\usepackage{amsmath}
\usepackage{graphicx}
\usepackage{subcaption}
\usepackage{booktabs}
\usepackage[section]{placeins}

\title{Implementacja i analiza sieci Hopfielda w zadaniu odtwarzania wzorców}
\author{Hubert Błonowski \and Jan Zaborski}
\date{}

\begin{document}

\maketitle

\begin{abstract}
Zaimplementowaliśmy od podstaw sieć Hopfielda z dwiema regułami uczenia: Hebba i Oji,
oraz dwoma trybami odtwarzania: synchronicznym i asynchronicznym.
Zbadaliśmy działanie sieci na zbiorach prostych bitmap, mierząc skuteczność metrykami
oraz szczegółowo analizując konkretne przypadki podatne na nieprawidłowe działanie sieci.
\end{abstract}

\tableofcontents

\section{Wstęp}
Rozważamy zadanie odtwarzania wzorców, polegające na odtworzeniu najbliższego zapamiętanego wzorca
ze zbioru wzorców zapamiętanych przez sieć na podstawie wzorca wejściowego.
Za wzorzec będziemy przyjmować wektor stanów.
Istotną różnicą między odtwarzaniem wzorców a klasyfikacją jest możliwość zwrócenia wzorca spoza zbioru wzorców zapamiętanych,
w szczególnych przypadkach przeciwnego wzorca (stany przeciwne do stanów zapamiętanego wzorca) lub mieszanki wzorców (część stanów z jednego, część z drugiego).
Do realizacji zadania wykorzystamy sieć Hopfielda.

Rozdział~\ref{sec:teoria} krótko omawia ideę sieci Hopfielda i ich uczenia,
rozdział~\ref{sec:impl} implementację,
a rozdział~\ref{sec:proc} procedurę eksperymentalną.
Wyniki przedstawia rozdział~\ref{sec:wyniki}.

\section{Podstawy teoretyczne}
\label{sec:teoria}

Sieć Hopfielda to sieć neuronowa, w której neurony $x_i$ przechowują stany
i każda para różnych neuronów jest połączona krawędzią o pewnej wadze (odpowiada to nieskierowanej, ważonej klice, gdzie każdy wierzchołek ma przypisaną liczbę).
Będziemy badać jedynie architekturę, w której stany należą do zbioru $\{-1,\, 1\}$, a wagi do zbioru liczb rzeczywistych.
Zatem wzorce zapamiętane, wejścia, wyjścia i stany pośrednie sieci będziemy oznaczać wektorem $\mathbf{x} \in \{-1,\, 1\}^{n}$,
gdzie $n$ to ustalona wspólna wymiarowość ich wszystkich (liczba neuronów).
Wagi będzie opisywać macierz $\mathbf{W} = \left[w_{ij}\right]_{i,\, j = 1, \ldots, n} \in \mathbb{R}^{n \times n}$,
gdzie $w_{ij}$ oznacza wagę krawędzi między $x_i$ a $x_j$.
Ponieważ krawędzie są nieskierowane, macierz jest symetryczna ($w_{ij} = w_{ji}$),
a $w_{ii} = 0$ odpowiada temu, że żaden neuron nie jest połączony z samym sobą.

\subsection{Uczenie}
Pamięć modelu będziemy kodować w wagach między neuronami.
Rozważamy dwie reguły wyznaczania macierzy $\mathbf{W}$: regułę Hebba (rozdz.~\ref{sec:hebb}) oraz regułę Oji (rozdz.~\ref{sec:oji}).

Oznaczamy $\mathbf{X} = \left[\mathbf{x}^{(p)} : p = 1, \ldots, N\right]$ jako zbiór wzorców zapamiętywanych przez sieć,
gdzie $N$ to liczba wzorców, $\mathbf{x}^{(p)}$ to wektor stanów wzorca o indeksie $p$,
a $\mathbf{X} \in \{-1,\, 1\}^{N \times n}$ jest macierzą, której kolejne wiersze odpowiadają kolejnym wzorcom (ich wektorom stanów).

\subsubsection{Reguła Hebba}
\label{sec:hebb}
Macierz wag obliczamy ze wzoru
\[
  \mathbf{W} = \frac{1}{n}\left(\mathbf{X}^{T}\mathbf{X} - N\mathbf{I}_{n}\right).
\]
Każdy element przekątnej $\mathbf{X}^{T}\mathbf{X}$ jest równy $N$, więc odjęcie $N\mathbf{I}_{n}$ zeruje przekątną.
Dodatni czynnik $\frac{1}{n}$ nie zmienia znaków pól lokalnych, a zatem nie wpływa na przebieg odtwarzania;
utrzymuje jedynie wagi i energię w podobnej skali niezależnie od liczby neuronów.

\subsubsection{Reguła Oji}
\label{sec:oji}

\subsection{Odtwarzanie wzorca}
Rozważamy dwa tryby odtwarzania wzorca: synchroniczny (rozdz.~\ref{sec:sync}) oraz asynchroniczny (rozdz.~\ref{sec:async}).
Indeks iteracji zapisujemy w indeksie górnym w nawiasach kwadratowych:
$\mathbf{x}^{[t]}$ oznacza stan sieci po $t$ iteracjach, a $\mathbf{x}^{[0]}$ wzorzec wejściowy
(nawiasy okrągłe, jak w $\mathbf{x}^{(p)}$, oznaczają numer wzorca zapamiętanego).
W obu trybach stan neuronu jest wielokrotnie ustawiany zgodnie ze znakiem jego \emph{pola lokalnego},
czyli ważonej sumy stanów pozostałych neuronów:
\[
  h_i = \sum_{j=1}^{n} w_{ij}\, x_j ,
  \qquad
  x_i \leftarrow
  \begin{cases}
    1,   & h_i > 0, \\
    -1,  & h_i < 0, \\
    x_i, & h_i = 0.
  \end{cases}
\]
Stan $\mathbf{x}$, w którym żaden neuron nie zmienia się po aktualizacji, nazywamy \emph{stanem stabilnym}.
Wzorzec zapamiętany jest odtwarzany poprawnie tylko wtedy, gdy jest stanem stabilnym sieci.

Przydatną wielkością jest \emph{energia} stanu
\[
  E(\mathbf{x}) = -\frac{1}{2}\, \mathbf{x}^{T}\mathbf{W}\mathbf{x}.
\]
Wzorce zapamiętane leżą zwykle w minimach lokalnych energii, a odtwarzanie można traktować jako schodzenie do najbliższego takiego minimum.

\subsubsection{Odtwarzanie synchroniczne}
\label{sec:sync}
W trybie synchronicznym wszystkie neurony są aktualizowane jednocześnie, na podstawie stanu z poprzedniej iteracji:
\begin{enumerate}
  \item Przyjmij $\mathbf{x}^{[0]}$ równe wzorcowi wejściowemu.
  \item W iteracji $t = 0, 1, \ldots$ oblicz pola lokalne $\mathbf{h}^{[t]} = \mathbf{W}\mathbf{x}^{[t]}$
        i dla wszystkich neuronów naraz ustaw
        \[
          x_i^{[t+1]} =
          \begin{cases}
            1,         & h_i^{[t]} > 0, \\
            -1,        & h_i^{[t]} < 0, \\
            x_i^{[t]}, & h_i^{[t]} = 0.
          \end{cases}
        \]
  \item Zakończ, gdy $\mathbf{x}^{[t+1]} = \mathbf{x}^{[t]}$ (stan stabilny),
        gdy $\mathbf{x}^{[t+1]} = \mathbf{x}^{[t-1]}$ (cykl długości 2)
        lub po osiągnięciu maksymalnej liczby iteracji.
\end{enumerate}
Dla symetrycznej macierzy wag odtwarzanie synchroniczne zawsze kończy się w stanie stabilnym
albo w cyklu długości 2, w którym sieć oscyluje między dwoma stanami~\cite{bruck1990}.

\subsubsection{Odtwarzanie asynchroniczne}
\label{sec:async}
W trybie asynchronicznym neurony są aktualizowane pojedynczo, a każdy od razu korzysta z aktualnych stanów pozostałych.
Iteracją jest tu jeden przebieg po wszystkich neuronach, a $\mathbf{x}^{[t]}$ oznacza stan po $t$ przebiegach:
\begin{enumerate}
  \item Przyjmij $\mathbf{x}^{[0]}$ równe wzorcowi wejściowemu.
  \item W przebiegu $t = 0, 1, \ldots$ zacznij od $\mathbf{x} = \mathbf{x}^{[t]}$
        i odwiedź wszystkie neurony w losowej kolejności;
        dla odwiedzanego neuronu $i$ oblicz $h_i$ na podstawie bieżącego stanu $\mathbf{x}$
        (z uwzględnieniem zmian wprowadzonych wcześniej w tym przebiegu) i ustaw
        \[
          x_i \leftarrow
          \begin{cases}
            1,   & h_i > 0, \\
            -1,  & h_i < 0, \\
            x_i, & h_i = 0.
          \end{cases}
        \]
        Stan po odwiedzeniu wszystkich neuronów to $\mathbf{x}^{[t+1]}$.
  \item Zakończ, gdy $\mathbf{x}^{[t+1]} = \mathbf{x}^{[t]}$, czyli w całym przebiegu żaden neuron nie zmienił stanu (stan stabilny),
        lub po osiągnięciu maksymalnej liczby przebiegów.
\end{enumerate}
Dla symetrycznej macierzy wag o zerowej przekątnej każda zmiana stanu pojedynczego neuronu obniża energię $E(\mathbf{x})$,
a stanów jest skończenie wiele, więc odtwarzanie asynchroniczne zawsze kończy się w stanie stabilnym~\cite{hopfield1982};
w przeciwieństwie do trybu synchronicznego nie występują tu oscylacje.

\section{Implementacja sieci}
\label{sec:impl}
Sieć zaimplementowano w języku C++ w standardzie C++26, korzystając jedynie z biblioteki standardowej.
Własnoręcznie zaimplementowano model macierzy dwuwymiarowej (i wybrane operacje z algebry liniowej),
oraz model sieci Hopfielda z API do uczenia sieci oraz generowania kolejnych kroków odtwarzania według wcześniej wspomnianych trybów.

\section{Procedura eksperymentalna}
\label{sec:proc}

\subsection{Zbiory danych}
\label{sec:dane}
Przebadano działanie sieci, wykorzystując 8 zbiorów danych jako pamięć modelu (tab.~\ref{tab:zbiory}).
Każdy zbiór to bitmapy o ustalonym rozmiarze $W \times H$ (szerokość $\times$ wysokość),
spłaszczone wierszami do wektorów stanów o długości $n = WH$;
piksel czarny odpowiada stanowi $1$, biały stanowi $-1$.

\begin{table}[htbp]
\centering
\begin{tabular}{lccc}
\toprule
zbiór & rozmiar $W \times H$ & neurony $n$ & wzorce $N$ \\
\midrule
\texttt{small-7x7}         & $7 \times 7$   & 49   & 5  \\
\texttt{letters-abc-8x12}  & $8 \times 12$  & 96   & 3  \\
\texttt{animals-14x9}      & $14 \times 9$  & 126  & 6  \\
\texttt{letters-14x20}     & $14 \times 20$ & 280  & 26 \\
\texttt{OCRA-12x30-cut}    & $12 \times 30$ & 360  & 7  \\
\texttt{large-25x25}       & $25 \times 25$ & 625  & 6  \\
\texttt{large-25x25.plus}  & $25 \times 25$ & 625  & 15 \\
\texttt{large-25x50}       & $25 \times 50$ & 1250 & 6  \\
\bottomrule
\end{tabular}
\caption{Zbiory wzorców zapamiętywanych przez sieć.}
\label{tab:zbiory}
\end{table}

Jako wzorce wejściowe do odtwarzania posłużyły zmodyfikowane wzorce z pamięci modelu:
\begin{itemize}
    \item \textbf{odwrócenie pikseli} -- $pn$ (po zaokrągleniu do liczby całkowitej) losowo wybranych pikseli
        zmienia stan na przeciwny, dla $p$ równego 5, 10, 15, 20, 25, 30, 40 i 50\%;
  \item \textbf{zaszumiony prostokąt} -- prostokąt o proporcjach obrazu i polu ok.\ $p$ obrazu, w losowym położeniu,
        zastąpiony losowymi stanami, dla $p$ równego 10, 20, 30, 40 i 50\%;
  \item \textbf{wymazany pas} -- pas o polu $p$ obrazu przy losowo wybranej krawędzi (górnej, dolnej, lewej lub prawej)
        ustawiony na $-1$ (tło), dla $p$ równego 25, 50 i 75\%;
  \item \textbf{przesunięcie} -- obraz przesunięty o 1, 2 lub 3 piksele w losowym kierunku,
        odsłonięty pas wypełniony stanem $-1$;
  \item \textbf{mieszanka} -- pierwsze $q$ kolumn pochodzi z wzorca źródłowego, pozostałe z innego, losowo wybranego wzorca
        tego samego zbioru, dla $q$ równego 50 i 75\%;
  \item \textbf{negatyw} -- wzorzec o wszystkich stanach przeciwnych, $-\mathbf{x}^{(p)}$.
\end{itemize}
Każdy wzorzec zapamiętany był źródłem 4 uszkodzonych wejść dla każdego rodzaju i poziomu uszkodzenia
(negatyw: jednego). Dodatkowo do sieci podawano 50 wejść losowych, w których każdy stan wynosi $1$ lub $-1$
z prawdopodobieństwem $\frac{1}{2}$, niezwiązanych z żadnym wzorcem.
Wszystkie wejścia generowano dla 5 ziaren generatora liczb losowych (0--4);
to samo ziarno wyznacza też kolejność aktualizacji neuronów w trybie asynchronicznym.
Odtwarzanie synchroniczne wzorców niezmienionych i negatywów jest deterministyczne, więc wykonano je jednokrotnie.

Ponieważ sieć Hopfielda ocenia wejście wyłącznie przez liczbę stanów różniących je od wzorców,
dla każdego wejścia mierzymy też \emph{odległość Hamminga} od wzorca źródłowego, czyli liczbę różniących się stanów.
Pozwala to porównywać różne rodzaje uszkodzeń na wspólnej skali.

\subsubsection{Zbiór \texttt{animals}}
\label{sec:ds_animals}
Zbiór składa się z 6 obrazków przedstawiających uproszczone sylwetki zwierząt.
Obrazki są rozmiaru $14 \times 9$~px.
Rysunek~\ref{fig:animals-14x9:patterns} przedstawia obrazki zawarte w zbiorze.
Rysunek~\ref{fig:animals-14x9:noise} przedstawia przykładowe modyfikacje wykorzystywane w eksperymentach.
Reszta zbiorów korzysta z tych samych modyfikacji, dlatego pominiemy dalsze przykłady.

\begin{figure}[htbp]
  \centering
  \includegraphics[width=0.72\linewidth]{figures/animals-14x9/patterns.png}
  \caption{Zbiór \texttt{animals}: 6 wzorców o rozmiarze $14\times9$ ($n = 126$). Czarny piksel to stan $1$, biały $-1$.}
  \label{fig:animals-14x9:patterns}
\end{figure}

\begin{figure}[htbp]
  \centering
  \includegraphics[width=\linewidth]{figures/animals-14x9/noise_examples.png}
  \caption{Przykłady uszkodzonych wejść utworzonych z wzorca 0 zbioru \texttt{animals}: wszystkie rodzaje i poziomy uszkodzenia.}
  \label{fig:animals-14x9:noise}
\end{figure}

\subsubsection{Zbiór \texttt{large}}
\label{sec:ds_large}
Zbiór składa się z 6 obrazków bez wspólnego motywu.
Obrazki są rozmiaru $25 \times 25$~px.
Rysunek~\ref{fig:large-25x25:patterns} przedstawia obrazki zawarte w zbiorze.

\begin{figure}[htbp]
  \centering
  \includegraphics[width=0.72\linewidth]{figures/large-25x25/patterns.png}
  \caption{Zbiór \texttt{large}: 6 wzorców o rozmiarze $25\times25$ ($n = 625$). Czarny piksel to stan $1$, biały $-1$.}
  \label{fig:large-25x25:patterns}
\end{figure}

\subsubsection{Zbiór \texttt{large plus}}
\label{sec:ds_large_plus}
Zbiór składa się z 15 obrazków, z czego 6 pochodzi ze zbioru \texttt{large} (rozdz.~\ref{sec:ds_large}).
Obrazki są tego samego formatu.
Rysunek~\ref{fig:large-25x25.plus:patterns} przedstawia obrazki zawarte w zbiorze.

\begin{figure}[htbp]
  \centering
  \includegraphics[width=0.96\linewidth]{figures/large-25x25.plus/patterns.png}
  \caption{Zbiór \texttt{large plus}: 15 wzorców o rozmiarze $25\times25$ ($n = 625$). Czarny piksel to stan $1$, biały $-1$.}
  \label{fig:large-25x25.plus:patterns}
\end{figure}

\subsubsection{Zmodyfikowany zbiór \texttt{large}}
\label{sec:ds_large_50}
Zbiór składa się z 6 obrazków, górna połowa obrazków to obrazki ze zbioru \texttt{large} (rozdz.~\ref{sec:ds_large}),
a dolna to jednolicie czarny kwadrat $25 \times 25$~px (wszystkie stany równe $1$).
Obrazki są rozmiaru $25 \times 50$~px.
Rysunek~\ref{fig:large-25x50:patterns} przedstawia obrazki zawarte w zbiorze.

\begin{figure}[htbp]
  \centering
  \includegraphics[width=0.72\linewidth]{figures/large-25x50/patterns.png}
  \caption{Zmodyfikowany zbiór \texttt{large}: 6 wzorców o rozmiarze $25\times50$ ($n = 1250$). Czarny piksel to stan $1$, biały $-1$.}
  \label{fig:large-25x50:patterns}
\end{figure}

\subsubsection{Zbiór \texttt{letters}}
\label{sec:ds_letters}
Zbiór składa się z 26 obrazków przedstawiających małe litery alfabetu.
Obrazki są rozmiaru $14 \times 20$~px.
Rysunek~\ref{fig:letters-14x20:patterns} przedstawia obrazki zawarte w zbiorze.

\begin{figure}[htbp]
  \centering
  \includegraphics[width=0.96\linewidth]{figures/letters-14x20/patterns.png}
  \caption{Zbiór \texttt{letters}: 26 wzorców o rozmiarze $14\times20$ ($n = 280$). Czarny piksel to stan $1$, biały $-1$.}
  \label{fig:letters-14x20:patterns}
\end{figure}

\subsubsection{Zbiór \texttt{letters-abc}}
\label{sec:ds_letters_abc}
Zbiór składa się z 3 obrazków przedstawiających 3 pierwsze małe litery alfabetu.
Litery mają inne kształty niż w zbiorze \texttt{letters} (rozdz.~\ref{sec:ds_letters}).
Obrazki są rozmiaru $8 \times 12$~px.
Rysunek~\ref{fig:letters-abc-8x12:patterns} przedstawia obrazki zawarte w zbiorze.

\begin{figure}[htbp]
  \centering
  \includegraphics[width=0.36\linewidth]{figures/letters-abc-8x12/patterns.png}
  \caption{Zbiór \texttt{letters-abc}: 3 wzorce o rozmiarze $8\times12$ ($n = 96$). Czarny piksel to stan $1$, biały $-1$.}
  \label{fig:letters-abc-8x12:patterns}
\end{figure}

\subsubsection{Zbiór \texttt{OCRA}}
\label{sec:ds_ocra}
Zbiór składa się z 7 obrazków przedstawiających wielkie litery alfabetu (A, B, C, K, Q, R, Z).
Obrazki są rozmiaru $12 \times 30$~px.
Rysunek~\ref{fig:OCRA-12x30-cut:patterns} przedstawia obrazki zawarte w zbiorze.

\begin{figure}[htbp]
  \centering
  \includegraphics[width=0.84\linewidth]{figures/OCRA-12x30-cut/patterns.png}
  \caption{Zbiór \texttt{OCRA}: 7 wzorców o rozmiarze $12\times30$ ($n = 360$). Czarny piksel to stan $1$, biały $-1$.}
  \label{fig:OCRA-12x30-cut:patterns}
\end{figure}

\subsubsection{Zbiór \texttt{small}}
\label{sec:ds_small}
Zbiór składa się z 5 obrazków przedstawiających proste kształty.
Obrazki są rozmiaru $7 \times 7$~px.
Rysunek~\ref{fig:small-7x7:patterns} przedstawia obrazki zawarte w zbiorze.

\begin{figure}[htbp]
  \centering
  \includegraphics[width=0.60\linewidth]{figures/small-7x7/patterns.png}
  \caption{Zbiór \texttt{small}: 5 wzorców o rozmiarze $7\times7$ ($n = 49$). Czarny piksel to stan $1$, biały $-1$.}
  \label{fig:small-7x7:patterns}
\end{figure}\section{Wyniki}
\label{sec:wyniki}

\section{Podsumowanie i wnioski}

\begin{thebibliography}{9}
  \addcontentsline{toc}{section}{Literatura}

\bibitem{hopfield1982}
J.~J.~Hopfield, \emph{Neural networks and physical systems with emergent collective computational abilities},
Proceedings of the National Academy of Sciences 79(8), 1982.

\bibitem{bruck1990}
J.~Bruck, \emph{On the convergence properties of the Hopfield model},
Proceedings of the IEEE 78(10), 1990.

\end{thebibliography}

\end{document}
